% GENERATED FILE. Edit canonical lessons, then run npm run generate:books.
% canonical-source-hash: fnv1a64:a6e8592c8a239913
% canonical-lessons: GU-C110-peti, GU-C110-barni, GU-C110-minbatti, GU-C110-takiyo, GU-C110-dhigli

\chapter{Box, Jar, Candle, Pillow, Doll}
\label{ch:gu-box-jar-candle-pillow-doll}
\hlchaptermodality{110}

\begin{chapteropening}
\textbf{By the end of this chapter:} \emph{I can say a box, a jar, a candle, a pillow and a doll.}

Complete the last lesson of chapter 110: i can say a box, a jar, a candle, a pillow and a doll.
\end{chapteropening}

\section[peṭī]{\gu{પેટી} (peṭī) — a box}

\label{lesson:GU-C110-peti}



\begin{quote}
Before the new one: say the Gujarati for rubbish, then the Gujarati for a basket.
\end{quote}

\subsection*{You'll want to know: \gu{પેટી}}
\textbf{\gu{પેટી}} (\emph{peṭī}) — \textquotedblleft{}a box\textquotedblright{}.

\begin{grammarlens}[title={What you've built}]
One more word for this chapter.
\end{grammarlens}

\subsection*{Your turn}
\begin{itemize}
\raggedright
  \item \emph{Say it:} \emph{peṭī}
  \item \emph{Say it:} \emph{peṭī}, once more
  \item \emph{Say it:} say \emph{peṭī}
  \item \emph{Recall:} say the Gujarati for rubbish, then the Gujarati for a basket, then say \emph{peṭī} again
\end{itemize}

\begin{tcolorbox}[breakable,colback=teal!4,colframe=teal!35!black,title={Before you move on}]
What does \gu{પેટી} mean? (\textquotedblleft{}A box\textquotedblright{}.) Say it once more.
\end{tcolorbox}
\section[barṇī]{\gu{બરણી} (barṇī) — a jar}

\label{lesson:GU-C110-barni}



\begin{quote}
Before the new one: say the Gujarati for a basket, then the Gujarati for a box.
\end{quote}

\subsection*{You'll want to know: \gu{બરણી}}
\textbf{\gu{બરણી}} (\emph{barṇī}) — \textquotedblleft{}a jar\textquotedblright{}.

\begin{grammarlens}[title={What you've built}]
One more word for this chapter.
\end{grammarlens}

\subsection*{Your turn}
\begin{itemize}
\raggedright
  \item \emph{Say it:} \emph{barṇī}
  \item \emph{Say it:} \emph{barṇī}, once more
  \item \emph{Say it:} say \emph{barṇī}
  \item \emph{Recall:} say the Gujarati for a basket, then the Gujarati for a box, then say \emph{barṇī} again
\end{itemize}

\begin{tcolorbox}[breakable,colback=teal!4,colframe=teal!35!black,title={Before you move on}]
What does \gu{બરણી} mean? (\textquotedblleft{}A jar\textquotedblright{}.) Say it once more.
\end{tcolorbox}
\section[mīṇbattī]{\gu{મીણબત્તી} (mīṇbattī) — a candle}

\label{lesson:GU-C110-minbatti}



\begin{quote}
Before the new one: say the Gujarati for a box, then the Gujarati for a jar.
\end{quote}

\subsection*{You'll want to know: \gu{મીણબત્તી}}
\textbf{\gu{મીણબત્તી}} (\emph{mīṇbattī}) — \textquotedblleft{}a candle\textquotedblright{}.

\begin{grammarlens}[title={What you've built}]
One more word for this chapter.
\end{grammarlens}

\subsection*{Your turn}
\begin{itemize}
\raggedright
  \item \emph{Say it:} \emph{mīṇbattī}
  \item \emph{Say it:} \emph{mīṇbattī}, once more
  \item \emph{Say it:} say \emph{mīṇbattī}
  \item \emph{Recall:} say the Gujarati for a box, then the Gujarati for a jar, then say \emph{mīṇbattī} again
\end{itemize}

\begin{tcolorbox}[breakable,colback=teal!4,colframe=teal!35!black,title={Before you move on}]
What does \gu{મીણબત્તી} mean? (\textquotedblleft{}A candle\textquotedblright{}.) Say it once more.
\end{tcolorbox}
\section[takiyo]{\gu{તકિયો} (takiyo) — a pillow}

\label{lesson:GU-C110-takiyo}



\begin{quote}
Before the new one: say the Gujarati for a jar, then the Gujarati for a candle.
\end{quote}

\subsection*{You'll want to know: \gu{તકિયો}}
\textbf{\gu{તકિયો}} (\emph{takiyo}) — \textquotedblleft{}a pillow\textquotedblright{}.

\begin{grammarlens}[title={What you've built}]
One more word for this chapter.
\end{grammarlens}

\subsection*{Your turn}
\begin{itemize}
\raggedright
  \item \emph{Say it:} \emph{takiyo}
  \item \emph{Say it:} \emph{takiyo}, once more
  \item \emph{Say it:} say \emph{takiyo}
  \item \emph{Recall:} say the Gujarati for a jar, then the Gujarati for a candle, then say \emph{takiyo} again
\end{itemize}

\begin{tcolorbox}[breakable,colback=teal!4,colframe=teal!35!black,title={Before you move on}]
What does \gu{તકિયો} mean? (\textquotedblleft{}A pillow\textquotedblright{}.) Say it once more.
\end{tcolorbox}
\section[ḍhī̃glī]{\gu{ઢીંગલી} (ḍhī̃glī) — a doll}

\label{lesson:GU-C110-dhigli}



\begin{quote}
Before the new one: say the Gujarati for a candle, then the Gujarati for a pillow.
\end{quote}

\subsection*{You'll want to know: \gu{ઢીંગલી}}
\textbf{\gu{ઢીંગલી}} (\emph{ḍhī̃glī}) — \textquotedblleft{}a doll\textquotedblright{}.

\begin{grammarlens}[title={What you've built}]
One more word for this chapter.
\end{grammarlens}

\subsection*{Your turn}
\begin{itemize}
\raggedright
  \item \emph{Say it:} \emph{ḍhī̃glī}
  \item \emph{Say it:} \emph{ḍhī̃glī}, once more
  \item \emph{Say it:} say \emph{ḍhī̃glī}
  \item \emph{Recall:} say the Gujarati for a candle, then the Gujarati for a pillow, then say \emph{ḍhī̃glī} again
\end{itemize}

\begin{tcolorbox}[breakable,colback=teal!4,colframe=teal!35!black,title={Before you move on}]
What does \gu{ઢીંગલી} mean? (\textquotedblleft{}A doll\textquotedblright{}.) Say it once more.
\end{tcolorbox}
